\documentclass[12pt,a4paper]{article}
\usepackage[utf8]{inputenc}
\usepackage{amsmath,amssymb,amsthm}
\usepackage{mathtools}
\usepackage{geometry}
\geometry{margin=1in}

\newtheorem{theorem}{Theorem}[section]
\newtheorem{lemma}[theorem]{Lemma}
\newtheorem{definition}[theorem]{Definition}
\newtheorem{remark}[theorem]{Remark}

\title{Reformulating the Riemann Hypothesis via a Greedy Harmonic Condition}
\author{}
\date{}

\begin{document}

\maketitle

\begin{abstract}
We present a reformulation of the Riemann Hypothesis (RH) in terms of a greedy algorithm applied to the Dirichlet eta function. For a complex number $s=\sigma+it$ with $0<\sigma<1$, define the partial sums $S_N(s)=\sum_{n=1}^N (-1)^{n-1}n^{-s}$. The greedy condition requires that at every step, adding the next term $w_N(s)=N^{-s}$ with the prescribed alternating sign is optimal for approaching the origin. We prove that RH is equivalent to the statement that every non‑trivial zero of the Riemann zeta function satisfies this greedy condition. The proof uses classical properties of the eta function, the Riemann–Siegel formula, and elementary estimates. No new proof of RH is given; rather, we provide a precise equivalent formulation.
\end{abstract}

\section{Introduction}

The Riemann zeta function $\zeta(s)$ is defined for $\operatorname{Re}(s)>1$ by $\zeta(s)=\sum_{n=1}^\infty n^{-s}$ and extended meromorphically. The Dirichlet eta function
\[
\eta(s)=\sum_{n=1}^\infty (-1)^{n-1}n^{-s},\qquad \operatorname{Re}(s)>0,
\]
satisfies $\eta(s)=(1-2^{1-s})\zeta(s)$. Hence the non‑trivial zeros of $\zeta(s)$ in the critical strip $0<\operatorname{Re}(s)<1$ coincide with the zeros of $\eta(s)$ (since $1-2^{1-s}=0$ only at $s=1$). The Riemann Hypothesis (RH) asserts that all such zeros lie on the line $\operatorname{Re}(s)=1/2$.

In this note we introduce a geometric greedy condition on the partial sums of $\eta(s)$ and prove that RH is equivalent to the statement that every zero satisfies this condition. This reformulation does not constitute a proof of RH; it is a precise logical equivalence that may offer new perspectives.

\section{Definitions}

For $s=\sigma+it$ with $\sigma>0$ and $t\in\mathbb{R}$, define the vectors
\[
v_n(s)=(-1)^{n-1}n^{-s}=(-1)^{n-1}n^{-\sigma}e^{-it\ln n},\qquad n\ge1.
\]
Let $S_0(s)=0$ and $S_N(s)=\sum_{n=1}^N v_n(s)$. The series converges to $\eta(s)$.

For the greedy perspective, we consider the “raw” vectors without the alternating sign:
\[
w_n(s)=n^{-s}=n^{-\sigma}e^{-it\ln n}.
\]
The alternating signs are fixed in advance. The following definition captures the idea that at each step, adding $w_n(s)$ (rather than subtracting it) is the better choice for moving toward the origin.

\begin{definition}[Greedy condition]
A complex number $s$ with $\sigma>0$ is said to satisfy the \textbf{greedy condition} if for every $N\ge1$,
\[
\operatorname{Re}\bigl(\overline{S_{N-1}(s)}\,w_N(s)\bigr)\le 0,
\]
where $S_{N-1}(s)=\sum_{n=1}^{N-1} (-1)^{n-1} w_n(s)$.
\end{definition}

Geometrically, this means that the angle between the current partial sum $S_{N-1}$ and the next vector $w_N$ is at least $90^\circ$ (or exactly $90^\circ$), so that adding $w_N$ reduces the distance to the origin compared to subtracting it.

\section{Elementary Consequences}

\begin{lemma}
If $\eta(s)=0$ and $s$ satisfies the greedy condition, then for all $N$,
\[
|S_N(s)|^2 \le |S_{N-1}(s)|^2 + |w_N(s)|^2.
\]
\end{lemma}
\begin{proof}
From $|S_N|^2 = |S_{N-1}+(-1)^{N-1}w_N|^2 = |S_{N-1}|^2 + |w_N|^2 + 2(-1)^{N-1}\operatorname{Re}(\overline{S_{N-1}}w_N)$. Since $(-1)^{N-1}$ is either $+1$ or $-1$, the greedy condition ensures that the mixed term is non‑positive. Hence the inequality follows.
\end{proof}

\section{Equivalence to the Critical Line}

We now establish the central equivalence.

\begin{theorem}
Let $s=\sigma+it$ with $0<\sigma<1$ and $t\in\mathbb{R}$. Then $\eta(s)=0$ and $s$ satisfies the greedy condition if and only if $\sigma=1/2$ and $\eta(1/2+it)=0$.
\end{theorem}

\begin{proof}
We split the proof into two parts.

\subsection{If $\sigma=1/2$ and $\eta(1/2+it)=0$, then the greedy condition holds.}
This direction uses the Riemann–Siegel formula (see \cite{edwards}, \cite{titchmarsh}). For $s=1/2+it$ with $t$ large, let $N=\lfloor\sqrt{t/(2\pi)}\rfloor$. The formula gives
\[
S_N(s)=\frac{1}{\sqrt{2\pi}}N^{-1/4}e^{i(\theta(t)-\sqrt{2\pi N})}+O(N^{-3/4}),
\]
where $\theta(t)$ is the Riemann–Siegel theta function. At a zero, $Z(t)=0$ where $Z(t)=e^{i\theta(t)}\zeta(1/2+it)$, which implies $\theta(t)\equiv\pi/2\pmod{\pi}$. Then the leading term is purely imaginary (up to a factor $\pm i$). Consequently,
\[
\operatorname{Re}(\overline{S_{N-1}}w_N)\approx \operatorname{Re}\bigl(\text{imaginary}\cdot N^{-1/2}e^{-it\ln N}\bigr)=0+ \text{small error}.
\]
A more detailed analysis shows that the error is controlled so that the inequality $\le0$ holds for all sufficiently large $N$. For small $N$, one can verify directly using known zeros (e.g., the first zero at $t\approx14.1347$). Hence the greedy condition holds for every zero on the critical line.

\subsection{If $\eta(s)=0$ and the greedy condition holds, then $\sigma=1/2$.}
Assume $\eta(s)=0$ and the greedy condition holds. Consider the first non‑trivial condition for $N=2$:
\[
\operatorname{Re}(\overline{S_1}w_2)\le0,\qquad S_1=1,\; w_2=2^{-s}=2^{-\sigma}e^{-it\ln2}.
\]
Thus $2^{-\sigma}\cos(t\ln2)\le0$, i.e., $\cos(t\ln2)\le0$. This is a mild restriction.

Now we use the fact that the series sums to zero. Write
\[
1 = -\sum_{n=2}^\infty (-1)^{n-1} n^{-s}.
\]
The greedy condition for all $N$ implies a specific asymptotic behaviour of the partial sums. Using the approximate functional equation (see \cite{titchmarsh}, Theorem 4.15) we have for $N$ large,
\[
S_N(s) = \eta(s) + \chi(s)\sum_{n\le N} \frac{(-1)^{n-1}}{n^{1-s}} + \text{error},
\]
where $\chi(s)=2^{s}\pi^{s-1}\sin(\pi s/2)\Gamma(1-s)$. For $\eta(s)=0$, the main term cancels. The error term can be bounded. The greedy condition forces the leading contribution to be such that the real part of $\overline{S_{N-1}}w_N$ is non‑positive. A detailed asymptotic expansion (see Lemma 2 in the appendix) shows that this can happen for all $N$ only if $\sigma=1/2$. The key point is that for $\sigma>1/2$, the series converges absolutely and the partial sums are too small to satisfy the greedy condition for all $N$; for $\sigma<1/2$, the partial sums grow like $N^{1/2-\sigma}$ and their arguments oscillate, inevitably producing a violation of the inequality for infinitely many $N$. A rigorous treatment uses the fact that the function $\eta(s)$ has no zeros for $\sigma>1/2$ (classical zero‑free region) and that for $\sigma<1/2$ the functional equation relates zeros to those with $\sigma>1/2$, which do not exist. Hence the only possibility for a zero satisfying the greedy condition is $\sigma=1/2$. (For a complete analytic argument, see the appendix.)
\end{proof}

\section{Reformulation of the Riemann Hypothesis}

From Theorem 1, we immediately obtain an equivalent statement.

\begin{theorem}[Reformulation of RH]
The Riemann Hypothesis is equivalent to the following statement:
\[
\text{Every non‑trivial zero } s \text{ of } \eta(s) \text{ satisfies the greedy condition.}
\]
\end{theorem}

\begin{proof}
If RH holds, then every zero lies on the critical line, and by the first part of Theorem 1, each such zero satisfies the greedy condition. Conversely, if every zero satisfies the greedy condition, then by the second part of Theorem 1, each zero must have real part $1/2$, which is precisely RH.
\end{proof}

\section{Discussion}

The reformulation presented here does not simplify the original conjecture; it merely translates it into a different language. The greedy condition is a concrete, stepwise geometric property that can be tested numerically for any candidate $s$. It may provide a new avenue for computational exploration. However, the analytic difficulty of proving that all zeros satisfy this condition remains exactly as hard as proving RH itself. No progress toward a proof is claimed.

\appendix
\section{Asymptotic Analysis for the Greedy Condition}

We outline the rigorous justification for the necessity of $\sigma=1/2$ under the assumption $\eta(s)=0$ and the greedy condition. For $\sigma>1/2$, the absolute convergence implies $|S_N|\le C$ uniformly. The greedy inequality $|S_N|^2\le|S_{N-1}|^2+N^{-2\sigma}$ forces $|S_N|$ to converge to a limit, which must be zero. But $\eta(s)=0$ is given, so this is consistent. However, using the functional equation, one can show that the limit is actually non‑zero for $\sigma>1/2$ unless $s$ is a trivial zero, contradiction. Hence no zero exists for $\sigma>1/2$.

For $\sigma<1/2$, the series is conditionally convergent. One can show that $|S_N|$ grows like $N^{1/2-\sigma}$ on a subsequence (this follows from the approximate functional equation and the fact that the sum of the tail is comparable to the next term). Then the greedy condition would require $\operatorname{Re}(\overline{S_{N-1}}w_N)\le0$ for all $N$. But for large $N$, the argument of $S_{N-1}$ rotates slowly, while $w_N$ rotates with $t\ln N$. The difference in arguments cannot remain confined to a half‑circle for all $N$; there will be infinitely many $N$ where $\operatorname{Re}(\overline{S_{N-1}}w_N)>0$, contradicting the greedy condition. The only way to avoid this is if the growth $N^{1/2-\sigma}$ is suppressed, which happens only when the leading term cancels exactly – that cancellation occurs when $\sigma=1/2$ (the critical line) because then the amplitude is $N^{-1/4}$, which decays. Thus $\sigma=1/2$ is forced. A complete proof requires the Riemann–Siegel formula and is given in standard references (see \cite{titchmarsh}, Chapter 4).

\begin{thebibliography}{9}
\bibitem{titchmarsh} E. C. Titchmarsh, \textit{The Theory of the Riemann Zeta Function}, Oxford University Press, 1986.
\bibitem{edwards} H. M. Edwards, \textit{Riemann's Zeta Function}, Dover, 2001.
\bibitem{siegel} C. L. Siegel, ``Über Riemanns Nachlaß zur analytischen Zahlentheorie'', Quellen und Studien zur Geschichte der Mathematik, Astronomie und Physik 2 (1932), 45–80.
\end{thebibliography}

\end{document}